%!TeX root=MemoriaTFG.tex

\chapter{Discusión}
\label{cap:discusion}

\section{Significado de los resultados}
\label{sec:significado}

La lectura agregada de los experimentos no apunta a la existencia de
un único modelo vencedor, sino a que el rendimiento dermatológico está
condicionado por cuatro ejes simultáneos: el nivel ontológico al que
se evalúa, el régimen de datos etiquetados disponible, la modalidad de
entrada y la estrategia de adaptación empleada. Bajo esta lectura, el
cuello de botella deja de residir únicamente en la calidad del encoder
visual y se desplaza hacia un conjunto de factores complementarios ---la
taxonomía con la que se organizan las clases, el volumen y la calidad de
las etiquetas, el alineamiento texto-imagen de la rama contrastiva, la
formulación de los \emph{prompts}, la equidad entre subgrupos y la propia
estrategia de despliegue---. Las subsecciones siguientes interpretan cada
uno de estos ejes a partir de la evidencia empírica del
capítulo~\ref{cap:resultados}, y la sección~\ref{sec:implicaciones}
deriva sus consecuencias operativas.

\subsection{Ausencia de un método óptimo único: especialización por tarea}
\label{sec:disc-especializacion}

La síntesis transversal de los experimentos sobre DermapixelAI~1.0
(tabla~\ref{tab:dermapixel-extra-sintesis}) muestra que ningún
codificador ni clasificador único domina todos los escenarios, sino que
cada tarea favorece una combinación distinta. PanDerm Large encabeza la
mayor parte de los escenarios de transferencia visual; DermLIP v2 y
Dermapixel R0 destacan sobre la cardinalidad clínica intermedia L2; el
\emph{retrieval} $k$-NN con índice FAISS resulta competitivo sobre L3 por
efecto de la cola larga; SAM2.1 domina la segmentación; los
\emph{ensembles} aportan en el cribado de seguridad de melanoma; e incluso
un codificador generalista no médico (CLIP-L) puede liderar la AUROC en
algún nivel pese a carecer de preentrenamiento clínico. El corolario es
que la elección del modelo no se resuelve en abstracto, sino en función
del objetivo: clasificación general, subcategoría clínica, cola larga,
ranking probabilístico, segmentación o cribado. Las subsecciones que
siguen y el resto de §\ref{sec:significado} desarrollan la evidencia que
sostiene esta especialización por tarea.

El primer hallazgo a discutir es el peso del preentrenamiento de
dominio. Los resultados del sondeo lineal
(sección~\ref{sec:rep-panderm} y sección~\ref{sec:benchmark})
sostienen la hipótesis de que preentrenar sobre datos clínicos
aporta una ventaja empírica medible frente a encoders generalistas
adaptados: PanDerm Large supera a PanDerm Base en accuracy, AUROC y
F1 ponderada sobre los diez datasets externos contrastados, con una
mejora media de $+9{,}0$\,pp en BAcc (tabla~\ref{tab:lp-panderm}).
El detalle por dataset de esa tabla matiza el origen del beneficio:
es modesto sobre dermatoscopia estandarizada ---donde una
arquitectura convolucional supervisada sobre ImageNet-21K llega a
empatar o superar la AUROC--- y crece sobre fotografía clínica con
dispositivos heterogéneos (PAD-UFES-20, DDI). Los resultados
sugieren que la ventaja del preentrenamiento de dominio aumenta con
la distancia entre el dataset y los datasets web generalistas, lo
que la hace especialmente relevante en condiciones clínicas no
estandarizadas, como la atención primaria o la teledermatología.

SigLIP-Large SO400M, un modelo vision-lenguaje generalista,
confirma esta lectura desde el extremo opuesto: lidera accuracy y
AUROC sobre HAM10000 (tabla~\ref{tab:benchmark}) ---atribuible a la
sobrerrepresentación de dermatoscopia estandarizada en los $400$
millones de pares imagen-texto web de su preentrenamiento--- pero
esa superioridad desaparece sobre PAD-UFES-20 y DDI. Su rango de
utilidad queda así delimitado a las modalidades canónicas, no al
material clínico heterogéneo.

El beneficio del preentrenamiento de dominio no es uniforme entre
niveles ontológicos, y esta segunda lectura matiza la anterior. El
experimento de combinación de codificadores sobre DermapixelAI~1.0
(sección~\ref{sec:dermapixel-ensemble},
tabla~\ref{tab:dermapixel-ensemble}) sitúa a DermLIP v2 como la
mejor cabeza individual sobre L2, con una ventaja de $+6{,}7$\,pp de
AUROC sobre PanDerm Large. El hallazgo es coherente con la lectura
\emph{zero-shot} sobre el mismo dataset
(sección~\ref{sec:dermapixel-zs}), donde DermLIP v2 lidera también
la AUROC L2 sin entrenar cabeza alguna. Leídos en conjunto, ambos
resultados sugieren que el preentrenamiento contrastivo texto-imagen
sobre el millón de pares de Derm1M~\cite{derm1m} captura la
granularidad de subcategoría clínica con mayor fidelidad que el
preentrenamiento auto-supervisado puro de la familia PanDerm. Esta
observación matiza la jerarquía PanDerm Large $>$ PanDerm Base $>$
DermLIP v2 observada en la mayoría de datasets y niveles: el
codificador óptimo depende del nivel ontológico, y el
preentrenamiento multimodal aporta una señal complementaria sobre la
categorización clínica intermedia que no se reduce al beneficio de
las \emph{features} visuales aisladas.

El ajuste fino parcial sobre DermapixelAI~1.0 nivel L1
(sección~\ref{sec:dermapixel-ft-l1},
tabla~\ref{tab:dermapixel-ft-multilevel}) refuerza esta lectura
desde un ángulo metodológico. Con $N_{\mathrm{train}} = 874$, el FT
que descongela parcialmente el encoder PanDerm Large mejora de forma
sustancial la BAcc sobre el sondeo lineal congelado, pero no supera
al \emph{ensemble} lineal \texttt{avg\_large\_dermlip}, que lo aventaja
tanto en BAcc como en AUROC (tabla~\ref{tab:dermapixel-ensemble}).
El patrón es coherente con la literatura sobre régimen de bajo $N$:
una combinación lineal de codificadores fundacionales congelados
aprovecha la diversidad de sus representaciones preentrenadas sin
añadir los parámetros que el FT denso debe ajustar sobre un dataset
reducido. La implicación de diseño para Dermapixel R0
(sección~\ref{sec:linea-spanderm}) es directa: priorizar la
adaptación de bajo rango mediante LoRA y la combinación de
codificadores con cabezas ligeras, antes que el FT denso clásico.

La ejecución posterior de la cabeza supervisada L2 de Dermapixel R0
(sección~\ref{sec:dermapixel-spanderm-v0}) confirma
cuantitativamente esta implicación y aporta la cifra-ancla del
argumento. La adaptación LoRA sobre los dos últimos bloques de
PanDerm Large, con sólo $524\,K$ parámetros entrenables
($0{,}17\%$ del total) frente a los $25{,}2$\,M del FT denso parcial
---una razón aproximada de $48\times$---, alcanza la mejor BAcc L2
documentada sobre el dataset, superando tanto al sondeo lineal
congelado como a DermLIP v2 individual. La lectura cruzada con
DermLIP v2 es particularmente informativa: la AUROC L2 de
Dermapixel R0 empata estadísticamente con la de DermLIP v2 nativo
(tabla~\ref{tab:dermapixel-ft-multilevel}). Esto sugiere que el
techo discriminativo del dataset a nivel L2 se alcanza tanto por el
preentrenamiento contrastivo PubMedBERT-DermLIP nativo como por la
adaptación textual ligera vía LoRA de un encoder de mayor capacidad:
la elección del codificador base y la estrategia de adaptación
operan como dimensiones aproximadamente ortogonales sobre la
cardinalidad L2.

La extensión posterior del FT denso parcial a los niveles L2 y L3
(sección~\ref{sec:dermapixel-ft},
tabla~\ref{tab:dermapixel-ft-multilevel}) cierra el contraste
metodológico. Sobre L2, el FT denso con $25{,}2$\,M parámetros
entrenables no supera a Dermapixel R0 LoRA ($524\,K$ parámetros):
queda por debajo en BAcc y empata en AUROC pese a la diferencia de
parametrización de $48\times$, lo que refuerza la hipótesis de
sección~\ref{sec:disc-lp-vs-ft} ---bajo $N_{\mathrm{train}} = 874$,
la adaptación de bajo rango sobre un encoder fundacional congelado
domina al FT denso clásico sobre el mismo encoder---. Sobre L3, en
cambio, el FT denso no mejora al sondeo lineal: ningún protocolo de
ajuste denso supera el techo que impone la cola larga estructural
($3{,}9$ muestras por clase de media en \emph{train}) sobre el
dataset actual.

La extensión a codificadores generalistas no médicos
(sección~\ref{sec:dermapixel-extra-encoders}) refuerza esta
no-uniformidad. CLIP-L-336, sin material biomédico en su
preentrenamiento, encabeza la AUROC bajo sondeo lineal sobre L2 y L3,
lo que matiza la hipótesis de especialización de dominio sobre las
métricas de ranking probabilístico. La sustitución del clasificador
paramétrico por \emph{retrieval} $k$-NN con índice FAISS sobre los
mismos \emph{embeddings} PanDerm Large
(sección~\ref{sec:dermapixel-extra-faiss}) supera al sondeo lineal en
el régimen de cola larga severa L3. Ambas observaciones cierran la
lectura de esta subsección: la elección óptima debe condicionarse al nivel
ontológico objetivo y al criterio operativo predominante, no a un
codificador o clasificador presupuesto como universal.

\subsection{DermapixelAI 1.0 como contribución metodológica}
\label{sec:disc-dermapixel-contrib}

Buena parte de la lectura anterior se apoya en DermapixelAI~1.0, cuyo
papel en este trabajo conviene precisar: no es un dataset auxiliar de
validación, sino una contribución metodológica en sí misma. Se trata de
un dataset clínico construido en castellano, independiente del dataset de
preentrenamiento Derm1M y auditado por hash MD5 frente a él
(sección~\ref{sec:auditoria-resultados}), lo que permite evaluar
transferencia, \emph{zero-shot}, adaptación LoRA y \emph{retrieval} sin la
incertidumbre de solapamiento que afecta a otros conjuntos. Su
organización conforme a la ontología jerárquica L1/L2/L3 habilita la
comparación de modelos a varios niveles de granularidad sobre un mismo
material, y su disponibilidad documentada lo sitúa como punto de partida
para evaluar modelos fundacionales sobre dermatología clínica en español.
La interpretación debe acotarse: se trata de una primera versión de
tamaño limitado y con una cola larga marcada a nivel L3, por lo que las
cifras derivadas requieren la estimación robusta discutida en
sección~\ref{sec:disc-cv-robustez} y no admiten extrapolación directa a
la práctica asistencial.

\subsection{Eficiencia en etiquetas y régimen de saturación}
\label{sec:disc-eficiencia}

Si el preentrenamiento de dominio es tan determinante, cabe
preguntarse cuántas etiquetas hacen falta para explotarlo. La curva
de eficiencia en etiquetas (sección~\ref{sec:label-eff-results},
tabla~\ref{tab:label-eff}) responde con un fenómeno propio de
encoders preentrenados sobre datasets masivos del dominio: el
rendimiento sobre HAM10000 satura con muy pocos datos etiquetados.
Con $1\%$ del conjunto de entrenamiento ($82$ imágenes),
PanDerm Large alcanza $0{,}850$ accuracy y $0{,}918$ AUROC. Eso
equivale al $95{,}7\%$ y al $96{,}2\%$ del rendimiento con el
conjunto completo. Con $5\%$ ($410$ imágenes), la AUROC sube a
$0{,}932$, el $97{,}7\%$ del máximo, a sólo $2{,}2$\,pp de la
cifra con $100\%$. La BAcc escala más despacio y no alcanza su
máximo hasta el $50\%$ del conjunto ($0{,}589$), lo que es
consistente con la dependencia de las clases minoritarias respecto
al volumen total.

Sobre PAD-UFES-20 ($N_{\mathrm{train}} = 1\,493$ con el conjunto
completo) la saturación es aún más temprana: con $14$ imágenes
($1\%$) se alcanza $0{,}740$ accuracy y $0{,}907$ AUROC, el
$95{,}5\%$ del AUROC con el conjunto completo. La implicación
clínica es directa. En contextos con pocos datos etiquetados
---enfermedades raras, centros de bajo volumen, subgrupos poco
representados--- una cabeza lineal sobre PanDerm Large preentrenado
es una opción técnicamente defendible.

\subsection{Sondeo lineal frente a ajuste fino}
\label{sec:disc-lp-vs-ft}

La elección entre congelar el encoder o reentrenarlo no tiene una
respuesta única, y esta subsección discute de qué depende. La
comparación directa entre sondeo lineal y ajuste fino sobre cuatro
datasets (sección~\ref{sec:ft-results},
tabla~\ref{tab:ft-panderm}) revela un patrón asimétrico con
implicaciones para el diseño de sistemas clínicos. Sobre HAM10000 y
PAD-UFES-20, datasets de tamaño intermedio o alto
($N \in [1\,493, 8\,207]$), el ajuste fino mejora la BAcc en
$+47{,}1$\,pp y $+19{,}4$\,pp respectivamente. Esto confirma la
utilidad de la receta original de PanDerm ---\emph{weighted random
sampler}, Mixup, CutMix y \emph{layer decay}--- para aprender las
clases minoritarias bajo desbalance acusado.

El comportamiento se invierte cuando el dataset es pequeño. Sobre
DDI, con $N_{\mathrm{train}} = 427$, el sondeo lineal supera al
ajuste fino en accuracy ($+7{,}3$\,pp), BAcc ($+2{,}1$\,pp),
AUROC ($+14{,}5$\,pp) y F1 ponderada ($+5{,}6$\,pp). El
descenso de AUROC ($0{,}827 \to 0{,}682$) es especialmente
diagnóstico: el encoder ajustado pierde capacidad de ranking sobre
el \emph{test}, lo que es compatible con sobreajuste a la
distribución de entrenamiento pese a las augmentaciones empleadas.
Sobre MSKCC ($N_{\mathrm{train}} = 6\,189$), la mejora de BAcc es
marginal ($+3{,}0$\,pp) y la AUROC desciende ligeramente
($-2{,}7$\,pp).

Estos tres casos delimitan un régimen operativo cuantificable. Para
datasets con $N_{\mathrm{train}} \lesssim 500$ imágenes, el sondeo
lineal sobre PanDerm Large es la opción de referencia. Para
$N_{\mathrm{train}} \gtrsim 5\,000$, el ajuste fino con la receta
original es preferible, siempre que existan clases minoritarias
clínicamente relevantes que el sondeo lineal no logre discriminar.
Entre ambos regímenes, la decisión depende del desbalance de clases
y del coste computacional asumible.

La oposición no se reduce, sin embargo, a congelar frente a ajustar el
encoder completo: cabe una tercera vía intermedia. La evidencia sobre
DermapixelAI~1.0 a nivel L2 (sección~\ref{sec:dermapixel-ft})
indica que la adaptación de bajo rango mediante LoRA mejora la BAcc frente
al ajuste fino parcial denso entrenando del orden de $48$ veces menos
parámetros. Este patrón condiciona la estrategia recomendable bajo $N$
reducido: la adaptación eficiente es preferible al ajuste denso, no sólo
por coste computacional sino por menor exposición al sobreajuste. La
implicación de diseño es directa y se materializa en la decisión
arquitectónica de Dermapixel R0 (sección~\ref{sec:dermapixel-spanderm-v0}),
que adopta LoRA sobre los últimos bloques de PanDerm Large en lugar de un
ajuste fino denso sobre el mismo encoder.

\subsection{Segmentación promptable y adaptación de bajo coste}
\label{sec:disc-segmentacion}

La segmentación aporta el hallazgo de mayor magnitud absoluta del
trabajo, y conviene discutir tanto su alcance como su origen. Sobre
ISIC2018 (sección~\ref{sec:seg-results}, tabla~\ref{tab:seg}),
SAM2.1-Large con \emph{fine-tuning} de su decodificador alcanza
Dice $0{,}947$ en el conjunto de test. Supera en $+2{,}6$\,pp la
cifra publicada por Yan et~al.~\cite{panderm2025} ($0{,}921$) y
en $+5{,}3$\,pp al baseline CAE-seg con backbone PanDerm Large
entrenado desde cero ($0{,}894$). El ajuste mantiene congelado el
\emph{image\_encoder} y entrena sólo el decodificador de máscaras y
el \emph{promptable encoder}. Por eso los pesos actualizados se
limitan a $\sim 4{,}2$\,M sobre los $\sim 227$\,M de
SAM2.1-Large (menos del $2\%$), con un coste computacional bajo.

La transferencia fuera de dominio es prácticamente nula: el Dice apenas
varía sobre ISIC2017 y PH2 respecto a ISIC2018. Este patrón sugiere que
las representaciones del encoder visual congelado (Hiera-Large) codifican
primitivas geométricas y de contorno suficientemente generales, de modo
que el ajuste del decodificador de máscaras y del \emph{promptable
encoder} captura los matices del dominio dermatoscópico sin sobreajustar a
una fuente de adquisición concreta. La magnitud del efecto de la
adaptación de dominio, cuantificada en sección~\ref{sec:seg-results}
frente a SAM2 \emph{zero-shot}, confirma que el coste de adaptación se
concentra en el decodificador y no en el encoder.

La comparativa pareada de arquitecturas promptables
(sección~\ref{sec:seg-arquitecturas},
tabla~\ref{tab:seg-arquitecturas}) matiza, no obstante, dónde reside
ese valor. La generación de la arquitectura pesa más que el
preentrenamiento médico: SAM2.1-tiny generalista supera a MedSAM~v1
con preentrenamiento médico, y el tamaño de \emph{backbone} no
aporta en el régimen \emph{box-prompted}. La especialización
dermatológica del segmentador parece, por tanto, un factor
secundario frente al diseño del modelo. Esta asimetría contrasta con
el peso decisivo del preentrenamiento de dominio observado en
clasificación (sección~\ref{sec:disc-especializacion}) e invita a
tratar por separado ambas familias de tareas.

El control automático sin caja
(sección~\ref{sec:seg-control-automatico}) reordena, sin embargo,
esta jerarquía de factores. Si las cinco arquitecturas
\emph{promptables} caben en $\sim 1$\,pp pero retirar el \emph{prompt}
cuesta $7{,}5$\,pp, el determinante real del rendimiento no es el
segmentador sino la caja: la información de localización pesa entre
seis y siete veces más que la elección de modelo. La consecuencia de
diseño es directa ---el esfuerzo de ingeniería rinde más en el
detector que produce la caja que en sustituir un segmentador por
otro--- y conecta con la decisión de producto de ofrecer un modo
interactivo, en el que el clínico aporta la caja, junto a uno
automático. Para este último, el análisis de robustez
(sección~\ref{sec:seg-robustez}) deja una especificación poco
habitual pero accionable: dado que apretar la caja degrada el Dice
mucho más que aflojarla, el detector debe sesgarse hacia cajas
ligeramente holgadas. La validación cruzada de cinco pliegues sitúa
además el \emph{headline} en $0{,}9554\pm 0{,}0010$, lo que eleva la
cifra de segmentación de estimación puntual a estimación con
intervalo, a diferencia del resto de protocolos del trabajo.

Conviene, por último, no sobreinterpretar la magnitud absoluta. El
análisis sobre anotación múltiple
(sección~\ref{sec:seg-techo-anotador}) muestra que el acuerdo entre
dermatólogos sobre el mismo contorno es de apenas $0{,}728$ de Dice,
mientras que el modelo concuerda con el consenso en $0{,}915$. El
$0{,}9556$ se sitúa, por tanto, cerca del techo que impone la propia
ambigüedad del borde lesional: las diferencias de $\sim 1$\,pp entre
arquitecturas son inferiores al desacuerdo inter-anotador y el «error»
residual frente a una referencia única es, en buena medida, ruido de
etiqueta. La salvedad pertinente es que esa consistencia convive con
una ligera sobreconfianza en los bordes dudosos, donde el modelo traza
una frontera nítida que los humanos no comparten.

\subsection{Sensibilidad al tokenizador en zero-shot multimodal}
\label{sec:disc-zs}

El rendimiento zero-shot no depende sólo del modelo, sino de cómo se
le interroga. La evaluación sobre HAM10000
(sección~\ref{sec:zs-results}, tabla~\ref{tab:zs-ham}) revela una
sensibilidad acusada al tokenizador y a la formulación de los
\emph{prompts}. La AUROC oscila entre $0{,}366$ (DermLIP v1/v2 con
tokenizador PubMedBERT y \emph{prompts} genéricos) y $0{,}854$
(DermLIP original con tokenizador GPT-2/CLIP y \emph{prompts} A3 de
Derm1M). Es un rango de $48{,}8$\,pp para variaciones que no tocan
los pesos visuales, sino sólo la rama lingüística y la formulación
textual de las clases.

El valor AUROC $<0{,}5$ de la configuración inicial es diagnóstico de un
desalineamiento entre el espacio textual proyectado por PubMedBERT y el
espacio visual del encoder de DermLIP v1/v2, hasta el punto de invertir la
correlación respecto a la tarea de clasificación. El experimento de
modificaciones aisladas (sección~\ref{sec:zs-results}) permite atribuir la
recuperación del rendimiento a la combinación de dos factores ---el cambio
de tokenizador y el cambio al conjunto de \emph{prompts} A3---, sin que la
evidencia disponible permita aislar una causa única de forma concluyente,
puesto que ambos operan sobre la misma rama textual. La lectura no debe,
por tanto, formularse como causalidad absoluta de un único componente,
sino como el efecto conjunto del tokenizador, los \emph{prompts} y el
alineamiento texto-imagen resultante. La rama textual no es un componente
accesorio del modelo multimodal: puede determinar el rendimiento final
tanto como la representación visual. De aquí se sigue una implicación
práctica directa: todo despliegue \emph{zero-shot} debe validar
empíricamente, sobre un conjunto del dominio, no sólo el modelo sino
también su tokenizador y la formulación de los \emph{prompts}, porque la
transferencia desde la configuración publicada por los autores no está
garantizada.

La evaluación adicional sobre DermapixelAI~1.0
(sección~\ref{sec:dermapixel-zs}, tabla~\ref{tab:dermapixel-zs})
extiende esta lectura al idioma. DermLIP v2 con \emph{prompts} en
inglés alcanza Acc@1 $= 0{,}500$ sobre la categoría etiológica
L1, frente a $0{,}361$ con \emph{prompts} equivalentes en
castellano ($-13{,}9$\,pp); SigLIP-SO400M cae $-11{,}1$\,pp en
la misma comparación. BiomedCLIP, en cambio, mantiene Acc@1
$= 0{,}389$ en ambos idiomas. Esto sitúa la penalización por
castellano como dependiente del dataset textual de preentrenamiento.
La magnitud de la caída motiva, como línea futura inmediata,
construir modelos vision-lenguaje contrastivos con rama textual
entrenada sobre material clínico en castellano.

\subsection{Brecha entre modelos específicos y LLMs multimodales}
\label{sec:disc-llm}

La sensibilidad de la rama contrastiva al texto invitaba a
comprobar si los grandes modelos de lenguaje multimodales, con
dataset de preentrenamiento mucho mayores, cierran la distancia con
los especialistas. La comparación de paradigmas sobre HAM10000
(sección~\ref{sec:llm-results}, tabla~\ref{tab:llm-comparativa})
muestra lo contrario y fija un orden de magnitud entre familias
relevante para elegir arquitectura en escenarios clínicos. El
gradiente de accuracy entre el mejor especialista ajustado (PanDerm
Base con ajuste fino y \emph{test-time augmentation}) y el mejor LLM
multimodal \emph{zero-shot} (GPT-4o) supera los $43$\,pp; la
tabla~\ref{tab:llm-comparativa} desglosa la escala completa de
configuraciones intermedias.

Tres observaciones cualitativas, detalladas en
sección~\ref{sec:llm-results}, matizan esta brecha sin reducirla.
Primero, la adaptación con LoRA sobre MedGemma~27B mejora la accuracy
sobre los datasets vistos en entrenamiento pero no transfiere de forma
sistemática al resto, lo que aconseja reportar su rendimiento por dataset
y no agregado. Segundo, GPT-4o y GPT-4o-mini exhiben un sesgo de error
hacia clases minoritarias del \emph{prompt} ---con sobreproducción de
predicciones de dermatofibroma y melanoma muy por encima de su prevalencia
real---, lo que sugiere que responden al espacio léxico de la consigna
tanto como a la evidencia visual. Tercero, BLIP-2 e InstructBLIP con
\emph{Q-Former} preentrenado producen una accuracy incompatible con
cualquier uso clínico, atribuible a la falta de adaptación al vocabulario
dermatológico. Las tres observaciones convergen en que la distancia entre
especialistas y LLMs multimodales generalistas no se explica por capacidad
bruta del modelo, sino por su grado de adaptación al dominio.

La brecha tiene además una vertiente de coste. La evaluación con
APIs comerciales sobre los tres datasets cuesta unos \$4{,}07.
MedGemma~27B, en cambio, opera localmente en BF16 sobre los
$128$\,GB de memoria unificada de la DGX Spark sin cuantización,
con latencia media de $1{,}49$ segundos por imagen. La diferencia
operativa importa para un despliegue hospitalario: los modelos
cerrados añaden coste recurrente y dependencia de un proveedor
externo, sin mejora apreciable de rendimiento clínico.

\subsection{Equidad por fototipo cutáneo}
\label{sec:disc-equidad}

La generalización entre fototipos es una condición de seguridad
clínica, y los resultados la cuantifican sin ambigüedad. La
evaluación sobre Fitzpatrick17k completo
(sección~\ref{sec:equidad-results}, tablas~\ref{tab:equidad-fototipo}
y~\ref{tab:cross-eval}) reproduce un patrón conocido en la
literatura: los cinco encoders contrastados presentan \emph{gap}~I-VI
positivo en la formulación de tres clases de malignidad, con valores
entre $+0{,}124$ y $+0{,}306$\,pp de BAcc. Ese \emph{gap} no se
interpreta en aislamiento, sino junto a la BAcc global de cada
modelo. DermLIP v2 obtiene la mejor BAcc global ($0{,}721$) y la
mejor AUROC global ($0{,}909$), pero su \emph{gap}~I-VI es
$+0{,}281$\,pp: el rendimiento sobre FST~VI ($0{,}434$) cae
mucho respecto a FST~III ($0{,}765$). PanDerm Large presenta un
\emph{gap} menor ($+0{,}217$\,pp) con BAcc global ligeramente
inferior ($0{,}699$). BiomedCLIP exhibe el \emph{gap} mínimo
($+0{,}124$\,pp), pero la peor BAcc global ($0{,}590$).

Esta última asimetría ---\emph{gap} mínimo con BAcc global
inferior--- ilustra que la métrica de equidad aislada no es
decisional. Un modelo uniformemente malo entre subgrupos puede
mostrar baja desigualdad aritmética sin ser equitativo en sentido
sustantivo, porque su rendimiento es bajo en todos los fototipos por
igual. Por eso BAcc global y \emph{gap} deben leerse de forma
conjunta. El mejor compromiso entre ambas variables es PanDerm
Large, seguido de DermLIP v2 cuando se ponderan el rendimiento sobre
FST~I-IV y la AUROC global.

El experimento cruzado \emph{Light}~$\leftrightarrow$~\emph{Dark}
(tabla~\ref{tab:cross-eval}) añade una segunda lectura. La caída de
rendimiento al entrenar sólo con fototipos claros y evaluar sobre
fototipos oscuros (L$\to$D vs L$\to$L) es máxima en DINOv2
($\Delta = -0{,}132$\,pp BAcc) y mínima en BiomedCLIP
($\Delta = -0{,}063$\,pp). DermLIP v2 mantiene la menor caída
Dark$\to$Light vs Dark$\to$Dark ($\Delta = +0{,}016$\,pp), lo que
sugiere mayor invariancia de sus representaciones al fototipo. El
patrón es coherente con su dataset de preentrenamiento Derm1M,
derivado de fuentes web abiertas con representación variable de
fototipos.

\subsection{Estimación robusta del rendimiento sobre datasets de tamaño moderado}
\label{sec:disc-cv-robustez}

La fiabilidad de las cifras sobre un dataset pequeño depende del método de
estimación empleado. La validación cruzada de cinco \emph{folds} agrupada
por caso clínico sobre DermapixelAI~1.0
(sección~\ref{sec:dermapixel-results}, tabla~\ref{tab:dermapixel-cv})
ilustra un punto metodológico relevante para evaluar modelos
fundacionales sobre datasets de tamaño moderado: la diferencia entre la
cifra puntual del \emph{split} fijo y la media de los \emph{folds} es
sustancial sobre las métricas dependientes de umbral (BAcc) y casi nula
sobre la AUROC \emph{one-vs-rest}. Esta observación matiza la lectura de
toda cifra de BAcc puntual sobre un \emph{test} reducido. De aquí se siguen
dos lecturas operativas. \textit{Primero}, la cifra del \emph{split}
publicado conserva validez como punto operativo de referencia ---y como
base de la comparación cruzada con la evaluación \emph{zero-shot} sobre el
mismo \emph{test} (sección~\ref{sec:dermapixel-zs})---, pero requiere la
estimación cruzada para defender el rendimiento esperado del sistema.
\textit{Segundo}, sobre datasets clínicos de tamaño moderado la AUROC es la
métrica primaria recomendable para comparar configuraciones de
\emph{transfer learning} congelado, mientras que la BAcc y la
\emph{accuracy} deben acompañarse de su intervalo de confianza o de su
estimación cruzada por \emph{folds}.

Leídas en conjunto, las ocho subsecciones anteriores indican que las
ganancias de rendimiento clínicamente relevantes no proceden de sustituir
un encoder por otro, sino de acertar en un conjunto de decisiones de
diseño: el nivel ontológico al que se evalúa y se reporta, la estrategia de
adaptación (sondeo lineal, LoRA o ajuste denso según el régimen de datos),
la formulación y el tokenizador de la rama textual, el equilibrio entre
BAcc y AUROC en presencia de desbalance, la validación por subgrupos de
fototipo y la combinación selectiva de modelos. Estas decisiones, más que
la elección aislada de un modelo, determinan el rendimiento alcanzable y
constituyen el objeto de las implicaciones operativas que se desarrollan en
la sección~\ref{sec:implicaciones}.

\section{Comparación con la literatura}
\label{sec:lit}

\subsection{Reproducción del benchmark de PanDerm}
\label{sec:lit-panderm}

El primer contraste con la literatura es la propia reproducción del
modelo. Yan et~al.~\cite{panderm2025} reportan para PanDerm Large,
sobre los $28$ \emph{downstream tasks} de su benchmark interno,
una mejora media del $+10{,}2\%$ sobre el estado del arte previo.
Las cifras de sondeo lineal de la tabla~\ref{tab:lp-panderm} caen
dentro del margen que reportan los autores para los datasets
compartidos (HAM10000, PAD-UFES-20, Derm7pt) y reproducen
cualitativamente el comportamiento del modelo sobre arquitectura
aarch64 con CUDA~13.0 y BF16. La verificación numérica frente al
hardware x86\_64 canónico es, por tanto, satisfactoria dentro de
las tolerancias esperables por aritmética BF16 y por diferencias de
versiones de \texttt{cuDNN}. Esto documenta la portabilidad del
pipeline al chip Grace Hopper GB10.

La cifra agregada del paper también es coherente con la observada
aquí. La diferencia PanDerm Large vs PanDerm Base sobre los diez
datasets externos contrastados ($+4{,}1$\,pp accuracy,
$+9{,}0$\,pp BAcc) sitúa el $+9{,}0$\,pp de BAcc cerca del
$+10{,}2\%$ que los autores reportan sobre su benchmark interno.
La ventaja del modelo Large parece, por tanto, estable entre los
conjuntos públicos contrastados y los conjuntos internos del grupo
de Monash.

\subsection{Segmentación binaria de lesión}
\label{sec:lit-seg}

En segmentación, la cifra obtenida también supera la referencia
publicada y conviene explicar por qué. SAM2.1-Large con
\emph{fine-tuning} del decodificador sobre ISIC2018 ($0{,}947$
Dice) mejora en $+2{,}6$\,pp el resultado de
Yan et~al.~\cite{panderm2025} ($0{,}921$ Dice con $100$ épocas).
Tres factores explican plausiblemente la diferencia. (i)~La
arquitectura promptable de SAM2.1 con \emph{bounding box} derivada
de la máscara de \emph{ground truth} aporta un guiado espacial
fuerte que las arquitecturas no promptables no aprovechan. (ii)~El
encoder Hiera-Large, preentrenado sobre $\sim 11$ millones de
imágenes y mil millones de máscaras, codifica primitivas
geométricas más informativas que las que el backbone visual de
PanDerm aprende sobre dermatoscopia. (iii)~El ajuste del
decodificador introduce los parámetros justos para capturar las
particularidades del dominio sin sobreajuste, mientras que un
\emph{fine-tuning} completo sobre $50$ épocas podría requerir más
datos o regularización adicional para superar la referencia. La
lectura operativa es que la segmentación binaria de lesión sobre
dermatoscopia roza la saturación cuando se combinan un encoder
promptable preentrenado a gran escala y una adaptación de bajo coste
sobre ISIC2018.

\subsection{Comportamiento zero-shot}
\label{sec:lit-zs}

El comportamiento zero-shot coincide en parte con lo publicado,
pero aporta también un matiz nuevo. La AUROC máxima alcanzada aquí
sobre HAM10000 con DermLIP en formulación zero-shot ($0{,}854$)
encaja cualitativamente con los rangos del paper original de
Derm1M~\cite{derm1m} para su mejor configuración. En cambio, la
sensibilidad al tokenizador documentada en
sección~\ref{sec:disc-zs} no está cuantificada de forma
sistemática en la literatura previa sobre la familia DermLIP, lo que
constituye una aportación empírica de este trabajo. La diferencia de
$48{,}8$\,pp AUROC entre la peor y la mejor configuración sobre el
mismo encoder visual es de un orden de magnitud comparable al efecto
de cambiar de paradigma (LP vs ZS). Esto desplaza la discusión
metodológica desde el modelo hacia su configuración operativa.

\subsection{Disparidad por fototipo}
\label{sec:lit-equidad}

La disparidad por fototipo confirma la literatura y, a la vez, la
amplía. El sentido de la desigualdad ---peor rendimiento sobre
fototipos elevados--- coincide con los hallazgos previos sobre
Fitzpatrick17k~\cite{fitzpatrick17k} y con la motivación que dio
origen a la curación de DDI~\cite{ddi}. Este trabajo extiende esa
línea evaluando cinco encoders bajo un mismo protocolo, lo que
permite contrastes directos. La paradoja BiomedCLIP (\emph{gap}
mínimo asociado a la peor BAcc global) no está documentada de forma
explícita en la literatura sobre equidad dermatológica. Es una
observación sobre la métrica \emph{gap} aislada que aconseja
reportes complementarios: BAcc global, AUROC global y \emph{gap}
deben presentarse conjuntamente.

\section{Implicaciones operativas}
\label{sec:implicaciones}

\subsection{Arquitectura defendible para sistemas de apoyo al diagnóstico}
\label{sec:disc-arq}

Los cuatro bloques de evidencia anteriores convergen en una
arquitectura técnicamente defendible para sistemas de apoyo al
diagnóstico dermatológico construidos sobre los modelos
fundacionales disponibles. La estructura recomendable es la
siguiente. El encoder es PanDerm Large preentrenado, congelado por
defecto y ajustado sólo bajo demanda. La cabeza es una regresión
logística multi-clase (L-BFGS, $C=1{,}0$,
\texttt{max\_iter}~$=5\,000$) cuando el dataset de entrenamiento es
reducido ($N_{\mathrm{train}} \lesssim 500$), o una cabeza
ajustada con la receta de PanDerm Large (Mixup, CutMix,
\emph{weighted random sampler}, $50$ épocas) cuando
$N_{\mathrm{train}} \gtrsim 5\,000$ y existe desbalance acusado.
La segmentación se resuelve con SAM2.1-Large adaptado mediante
\emph{fine-tuning} del decodificador cuando la tarea requiera
localizar la lesión.

Las ramas textuales se incorporan con cautela. La rama lingüística
contrastiva (DermLIP) sólo entra con prompts validados y tokenizador
compatible, verificado sobre un conjunto de validación del dominio.
La rama LLM multimodal generalista (GPT-4o, Gemini, MedGemma) no se
usa como clasificador principal, dado el orden de magnitud de
rendimiento inferior cuantificado en sección~\ref{sec:disc-llm}.
Su rol natural es complementario: generar razonamiento textual o
explicaciones para casos ya clasificados por el encoder
especializado.

La consecuencia de diseño que se desprende del conjunto de hallazgos
(sección~\ref{sec:hallazgos}) es que la arquitectura defendible es
modular, no un modelo único: PanDerm Large para la transferencia visual
general, adaptación LoRA para la subcategoría clínica bajo $N$ reducido,
\emph{retrieval} FAISS para la cola larga, SAM2.1 para la segmentación y
un \emph{ensemble} acotado para el cribado de melanoma, seleccionando el
componente según la tarea y el nivel ontológico. Sobre esta estructura,
dos validaciones son condición necesaria antes del despliegue: la de los
\emph{prompts} y el tokenizador de cualquier rama contrastiva sobre un
conjunto del dominio, dada la sensibilidad documentada en
sección~\ref{sec:disc-zs}, y la validación por fototipo cutáneo que se
detalla en sección~\ref{sec:disc-ventana}.

\subsection{Armonización taxonómica y nivel de agregación}
\label{sec:disc-armonizacion}

La ontología no sólo permite comparar datasets, sino que obliga a
elegir a qué nivel se compara. La heterogeneidad taxonómica entre
datasets, mitigada mediante la ontología jerárquica L1/L2/L3 descrita
en sección~\ref{sec:ontologia}, sigue requiriendo una decisión
explícita sobre el nivel de agregación de los contrastes. La
comparación inter-dataset es defendible a nivel L1 (cuatro familias
etiológicas) y L2 ($43$ subcategorías clínicas) sobre el dataset
armonizado de $72\,654$ imágenes. El nivel L3 ($367$
diagnósticos individuales) conserva la cola larga propia de cada
dataset y no admite contrastes directos sin agregación adicional. La
implicación para un sistema de producción es que la jerarquía debe
mostrarse de forma explícita al usuario clínico: presentar la
predicción L3 junto a su agregación L2 y L1 da robustez frente a la
incertidumbre de la cola larga sin sacrificar especificidad cuando
ésta es fiable.

\subsection{Ventana operativa por fototipo cutáneo}
\label{sec:disc-ventana}

El rendimiento desigual entre fototipos también acota dónde puede
desplegarse el sistema con garantías. Esa ventana operativa, definida
en el experimento de equidad (tabla~\ref{tab:equidad-fototipo}),
condiciona la generalización a poblaciones con fototipos elevados.
Los cinco encoders presentan \emph{gap}~I-VI positivo entre
$+0{,}124$ y $+0{,}306$\,pp BAcc. La paradoja BiomedCLIP
---\emph{gap} mínimo asociado a la peor BAcc global--- vuelve a
ilustrar que el \emph{gap} aislado no es una métrica decisional. El
mejor compromiso entre BAcc global y \emph{gap} sobre la formulación
de tres clases es PanDerm Large, seguido de DermLIP v2 cuando se
prioriza la AUROC global. La implicación clínica es clara: desplegar
sobre poblaciones con representación significativa de fototipos
elevados exige validación específica por subgrupos y mecanismos de
revisión humana cuando el modelo opere en su régimen de menor
confianza.

\subsection{Ensemble de seguridad para detección de melanoma}
\label{sec:disc-ensemble}

La detección de melanoma admite un esquema orientado a la seguridad,
y conviene discutir tanto su valor como sus costes. El protocolo
auxiliar de ensemble (sección~\ref{sec:ensemble-results},
tabla~\ref{tab:ensemble}) muestra que la combinación OR
\emph{top-3} de PanDerm Large \emph{fine-tuned} sobre HAM10000 (M1)
y SigLIP-Large SO400M con sondeo lineal alcanza \emph{recall} de
melanoma del $100\%$ ($70$ de $70$) sobre el \emph{split} de
test de HAM10000. El análisis de errores explica por qué: los dos
melanomas que M1 no captura en \emph{top-3} sí los captura SigLIP,
lo que confirma la complementariedad de los \emph{backbones} y
justifica la elección del par. M7 (clasificador unificado sobre el
dataset armonizado de $43$ clases) no añade melanomas adicionales
sobre HAM, pero mantiene valor en escenarios multidataset por su
mayor cobertura ontológica.

Tres caveats acompañan la lectura operativa de este resultado.
(i)~El \emph{recall} del $100\%$ es \emph{top-3}, no
\emph{top-1}, por lo que exige una interfaz que presente al clínico
las tres clases más probables. (ii)~El tamaño muestral de melanomas
es $N = 70$, lo que produce intervalos de confianza al $95\%$
amplios y obliga a validación prospectiva sobre una cohorte mayor.
(iii)~El coste en falsos positivos del esquema OR \emph{top-3} es
elevado ($885$ FP sobre $1\,162$ no-melanomas): asumible en un
cribado oportunista, pero suficiente para descartar su uso aislado
como clasificación primaria. La lectura operativa, coherente con la
arquitectura modular de sección~\ref{sec:disc-arq}, es que este
\emph{ensemble} debe integrarse como mecanismo de apoyo orientado a la
sensibilidad ---priorizar la detección de melanoma asumiendo falsos
positivos--- y nunca como decisión diagnóstica autónoma.

\section{Hallazgos no anticipados}
\label{sec:hallazgos}

Seis observaciones transversales del trabajo no se anticipaban desde la
literatura y merecen atención específica. Se ordenan de mayor a menor
alcance interpretativo, comenzando por las que afectan a la elección de
método y terminando por las acotadas a un protocolo concreto.

\paragraph{Ausencia de un método óptimo único.}
Ningún codificador ni clasificador domina todos los niveles ontológicos y
todas las tareas: PanDerm Large encabeza la transferencia visual general,
DermLIP v2 y Dermapixel R0 la subcategoría clínica L2, el \emph{retrieval}
$k$-NN la cola larga L3, SAM2.1 la segmentación y los \emph{ensembles} el
cribado de melanoma, mientras que un codificador generalista no médico
puede liderar la AUROC en algún nivel
(tabla~\ref{tab:dermapixel-extra-sintesis}). La literatura sobre modelos
fundacionales dermatológicos tiende a buscar un modelo de referencia
único; la evidencia aquí sugiere que la pregunta correcta es qué modelo
para qué tarea y a qué nivel ontológico.

\paragraph{Sensibilidad extrema al tokenizador y al \emph{prompt}.}
La AUROC \emph{zero-shot} de la familia DermLIP sobre HAM10000 varía en un
rango comparable al de cambiar de paradigma completo, sin modificar los
pesos del encoder visual, en función del tokenizador y de la formulación
de los \emph{prompts} (sección~\ref{sec:disc-zs}). Esta sensibilidad no
está cuantificada de forma sistemática en la literatura previa sobre
DermLIP y desplaza parte de la atención metodológica del modelo hacia su
configuración operativa.

\paragraph{Sondeo lineal supera al ajuste fino en DDI.}
Sobre DDI ($N_{\mathrm{train}} = 427$), el sondeo lineal sobre PanDerm
Large supera al ajuste fino en accuracy, AUROC, BAcc y F1 ponderada
(sección~\ref{sec:disc-lp-vs-ft}). El comportamiento, atribuible a
sobreajuste bajo $N$ reducido, no se documenta de forma explícita en la
literatura sobre PanDerm y ofrece una heurística operativa para datasets
pequeños.

\paragraph{Retrieval competitivo en la cola larga L3.}
La sustitución del clasificador lineal por \emph{retrieval} $k$-NN con
índice FAISS sobre los mismos \emph{embeddings} PanDerm Large supera al
sondeo lineal en el régimen de cola larga severa L3 de DermapixelAI~1.0
(sección~\ref{sec:dermapixel-extra-faiss}). El hallazgo indica que, cuando
el número de muestras por clase es muy bajo, un esquema no paramétrico
basado en vecindad puede ser preferible al ajuste de una frontera de
decisión paramétrica.

\paragraph{Asimetría rendimiento-equidad en DermLIP v2.}
DermLIP v2 alcanza la mejor BAcc y AUROC globales sobre Fitzpatrick17k pero
uno de los \emph{gap}~I-VI más altos entre los cinco encoders contrastados
(sección~\ref{sec:disc-equidad}). Que un rendimiento agregado superior
conviva con una desigualdad por subgrupo elevada ---en contraste con la
paradoja BiomedCLIP, que asocia \emph{gap} mínimo a la peor BAcc global---
sugiere que el preentrenamiento sobre Derm1M produce representaciones
informativas sobre la categoría diagnóstica pero no suficientemente
invariantes al fototipo cutáneo.

\paragraph{Interpretabilidad emergente del SAE Large sobre
Derm7pt.}
La validación independiente del SAE Large ($16\,384$
\emph{features}) sobre el dataset Derm7pt~\cite{kawahara2019},
reportada en sección~\ref{sec:sae-derm7pt}, identifica
\emph{features} individuales con AUROC efectivo hasta $0{,}823$
sobre \emph{vascular structures} y hasta $0{,}926$ sobre vasos
en horquilla ($N_{+} = 15$ positivos). Las anotaciones del
\emph{Seven-Point Checklist} no se usaron ni en el preentrenamiento
del SAE ni en el del encoder PanDerm Large. Esto implica que la base
sparse aprendida sobre Derm1M ($507\,657$ imágenes) ya ha
capturado representaciones alineadas con la taxonomía dermatoscópica
clásica. El hallazgo sostiene la interpretabilidad mediante Sparse
Autoencoders más allá de los $48$ conceptos visuales de
SkinCon~\cite{skincon} sobre los que se diseñó el experimento
original del Anexo~\ref{cap:anexof}. El \emph{Concept Bottleneck
Model} con cuello de botella sobre los siete criterios cuantifica
además el coste de esta interpretabilidad sobre la tarea binaria
melanoma/no-melanoma ($-5{,}8$\,pp AUROC, $0{,}890 \to
0{,}832$). La jerarquía de pesos del meta-clasificador
(\emph{blue whitish veil} $+2{,}54$, \emph{regression structures}
$+2{,}07$) es coherente con el peso clínico de estos criterios en
la escala diagnóstica de referencia~\cite{kawahara2019}. La lectura
operativa es que el SAE no es sólo un mecanismo de reducción
dimensional: admite interpretación concepto-a-concepto sobre el
vocabulario dermatoscópico. Esto refuerza la viabilidad del módulo de
interpretabilidad del prototipo (Anexo~\ref{cap:anexoh}) sin
reentrenar el SAE.

Dos observaciones de menor alcance completan el cuadro. La
competitividad de ConvNeXt-Large preentrenado sobre ImageNet-21K sobre
HAM10000, que supera ligeramente a PanDerm Large en AUROC, es un caso
particular del primer hallazgo: confirma que la ventaja del
preentrenamiento de dominio se concentra en el material clínico no
estandarizado y se diluye sobre dermatoscopia canónica. Y la adaptación
LoRA sobre MedGemma~27B mejora la accuracy sobre los datasets vistos en
entrenamiento sin transferir de forma sistemática al resto
(sección~\ref{sec:disc-llm}), lo que aconseja reportar su rendimiento por
dataset; ambas refinan hallazgos ya enunciados más que constituir patrones
independientes.

\section{Limitaciones interpretativas}
\label{sec:disc-limitaciones}

Conviene cerrar la discusión acotando hasta dónde llegan estas
conclusiones. Cuatro elementos condicionan la interpretación de los
resultados, y el capítulo~\ref{cap:limitaciones} los desarrolla en
detalle. Primero, el solapamiento exacto de Dermnet con el dataset de
preentrenamiento Derm1M ($100\%$ por hash MD5) introduce una
incertidumbre acotada sobre las cifras de la familia DermLIP en ese
dataset ---HIBA y MSKCC, pese a su origen ISIC-archive, no presentan
solapamiento MD5---; las cifras de PanDerm, DINOv2, ConvNeXt-Large,
EfficientNetV2-Large y SigLIP-Large no se ven afectadas. Segundo, los
tamaños muestrales de los subgrupos Fitzpatrick V y VI ($169$ y
$73$ imágenes respectivamente en el conjunto de test) producen
intervalos de confianza amplios para las cifras de \emph{gap}, lo que
limita la cuantificación exacta de la desigualdad por fototipo.
Tercero, la sensibilidad al \emph{prompt} en zero-shot y en la
evaluación de LLMs multimodales hace que las cifras reportadas sean
condicionales a una formulación concreta; generalizarlas a otras
formulaciones requiere validación específica. Cuarto, el trabajo
evalúa rendimiento sobre fotografías de archivos de investigación; la
generalización a imagen real adquirida en flujos de teledermatología
sobre cohorte hospitalaria prospectiva no se aborda y constituye una
línea futura inmediata.

Estas limitaciones acotan el dominio de validez de las cifras, pero
no invalidan los patrones cualitativos identificados: la ventaja del
preentrenamiento de dominio, la saturación temprana del sondeo
lineal, el régimen $N_{\mathrm{train}} \lesssim 500$ como umbral de
preferencia del LP, la brecha entre especialistas y LLMs
multimodales generalistas, y la existencia de un \emph{gap} por
fototipo sistemático y consistente entre encoders.
